\documentclass[conference]{IEEEtran}
\usepackage{amsmath}
\usepackage{booktabs}
\usepackage{url}

\title{Patient-to-Criterion Eligibility Screening with a Resource-Constrained Transformer}

\author{\IEEEauthorblockN{Student Name}
\IEEEauthorblockA{Course: Generative AI\\
Institution: To be completed}}

\begin{document}
\maketitle

\begin{abstract}
This paper describes a reproducible prototype for matching patient descriptions to
individual clinical-trial eligibility criteria. The prescribed TrialGPT Criterion
Annotations dataset is partitioned by patient, and a Qwen2.5 decoder model is adapted
with parameter-efficient quantized fine-tuning. The system preserves expert label
categories and sentence-level evidence IDs, and implements reference and optimized
inference paths for key/value caching, tiled online-softmax attention, and greedy
speculative decoding. A browser application exposes criterion-level predictions and
inference comparisons. Experimental results are intentionally pending execution on a
CUDA runtime; no performance values are claimed in this draft.
\end{abstract}

\begin{IEEEkeywords}
clinical trial matching, eligibility criteria, evidence retrieval, QLoRA, key/value
cache, online softmax, speculative decoding
\end{IEEEkeywords}

\section{Introduction}
Clinical trial screening requires comparing patient facts with many individual
inclusion and exclusion criteria. TrialGPT established criterion-level matching with
expert annotations and evidence-oriented explanations \cite{jin2024trialgpt}. This
project implements a smaller local-model system under a 4\,GB RTX 3050 constraint,
while keeping training portable to free Colab or Kaggle GPU runtimes.

The unit of prediction is a patient--criterion pair. The system does not treat an
aggregated patient--trial verdict as a separately labeled ground truth. Missing
information is not treated as evidence of either eligibility or ineligibility.

\section{Related Work}
TrialGPT decomposes trial matching into retrieval, criterion matching, and ranking;
this work focuses on the labeled criterion-level task \cite{jin2024trialgpt}. LoRA
adapts frozen pretrained weights with low-rank trainable matrices
\cite{hu2022lora}; QLoRA combines this approach with 4-bit quantization to reduce
fine-tuning memory \cite{dettmers2023qlora}. FlashAttention motivates tiled,
IO-aware exact attention and online normalization \cite{dao2022flashattention}.
Speculative decoding verifies draft tokens with a target model while preserving
greedy outputs when the acceptance procedure is correct
\cite{leviathan2023speculative}.

\section{Materials and Methods}
\subsection{Dataset and Split}
The required TrialGPT Criterion Annotations release contains 1,015 criterion
annotations from 53 unique patient IDs \cite{trialgpt_dataset}. Patient IDs are sorted,
shuffled with NumPy \texttt{default\_rng(42)}, and allocated to 37 training, 8
validation, and 8 test patients. The resulting annotation counts are 603, 206, and
206, respectively. The generated manifest records patient IDs and counts. No test
labels are used for tuning.

The model input uses only the patient note, criterion type, and criterion text. The
expert eligibility label and expert sentence IDs are targets. Dataset fields
\texttt{gpt4\_explanation}, \texttt{gpt4\_sentences}, \texttt{gpt4\_eligibility}, and
\texttt{explanation\_correctness} are excluded. One source annotation has a null
criterion text; it is retained with an empty value and an explicit missing-field
marker rather than a synthesized criterion.

\subsection{Models and Fine-Tuning}
Two open-weight decoder sizes from the same model family were considered. The
selected target is Qwen2.5-1.5B-Instruct (1.54B parameters, 28 decoder layers, grouped
query attention, Apache-2.0); Qwen2.5-0.5B-Instruct (0.49B parameters, 24 layers,
Apache-2.0) is the tokenizer-compatible draft model
\cite{qwen25_15b_modelcard,qwen25_05b_modelcard}. At 4 bits, parameter payload alone
has arithmetic lower bounds of approximately 0.72\,GiB and 0.23\,GiB, respectively;
these are not peak-memory measurements and exclude quantization metadata, activations,
KV cache, and runtime overhead.

The training script uses 4-bit NF4 with double quantization, LoRA rank 8 and
$\alpha=16$, microbatch size 1, gradient accumulation, and gradient checkpointing.
The training target is a JSON object containing the original expert category and
expert evidence sentence IDs. The loss is applied to target tokens and not to prompt
tokens. The default maximum sequence length is 1,024 tokens with a reserved output
budget. Long notes are truncated by removing trailing numbered sentences first, then
by token prefix if necessary. Per-example metadata records truncation and any gold
evidence IDs omitted from the prompt.

\subsection{Criterion and Evidence Evaluation}
Classification is evaluated with accuracy, macro-F1, per-class precision, recall and
F1, and confusion matrices, overall and separately for inclusion and exclusion
criteria. Invalid JSON outputs are counted under an explicit invalid-output label.
Evidence retrieval uses set-based sentence IDs and reports micro and macro precision,
recall, and F1. For a record with no gold evidence, an empty prediction is an exact
empty match (per-record precision, recall, and F1 equal 1); a non-empty prediction is
scored as 0 for that record. Invalid or out-of-context citations are separately
counted. Explanations quote source patient sentences rather than generated
unverifiable text.

\subsection{KV-Cached Decoding}
The uncached reference recomputes the complete prefix at every greedy step. The cached
path performs a prompt prefill, then passes only the newest token and reuses the
returned key/value tensors. It supplies explicit cache positions and position IDs,
extends the attention mask as the cache grows, and verifies the cache length after
each update. Greedy token sequences are compared exactly. The benchmark records
prefill time, decode time, throughput, and incremental peak CUDA allocation.

\subsection{Tiled Online-Softmax Attention}
The reference computes scaled dot-product attention
\begin{equation}
  \mathrm{Attention}(Q,K,V)
  = \mathrm{softmax}\left(\frac{QK^\top}{\sqrt{d_k}} + M\right)V,
\end{equation}
where $M$ encodes causal and optional boolean masks. The tiled implementation updates
the running row maximum $m$, normalization $l$, and weighted value accumulator $o$
for each key block:
\begin{align}
  m' &= \max(m, m_b), \\
  \alpha &= \exp(m-m'), \quad
  p_b = \exp(s_b-m'), \\
  l' &= \alpha l + \sum p_b, \quad
  o' = \alpha o + p_b V_b.
\end{align}
The result is $o/l$, with all-masked rows defined as zero. The provided educational
implementation uses NumPy and is checked against the reference on controlled causal
and masked inputs; it is not a fused GPU kernel.

\subsection{Speculative Decoding}
The draft model proposes a short greedy token block. The target verifies the block in
one forward pass; matching tokens are accepted until the first mismatch, at which
point the target's greedy token is emitted. Both caches are cropped to the accepted
prefix and advanced with the correction token. If every draft token is accepted, a
target bonus token may be emitted. Tokenizer vocabularies are checked for exact
compatibility. Benchmarks report token equality, acceptance rate, latency, throughput,
and incremental memory; speedup is not assumed.

\subsection{Browser Application}
The Streamlit application accepts a patient record and criterion list, returns one
result per criterion, displays exact evidence sentences, and exports a Markdown
screening report. A prominent notice states that the prototype is educational and is
not a clinical enrollment decision. The interface also exposes KV-cache and
speculative-decoding comparisons.

\section{Experimental Setup and Metrics}
The local development machine is Windows with an NVIDIA GeForce RTX 3050 and 4\,GB
VRAM. PyTorch/CUDA was not installed in the development Python environment during
initial implementation. QLoRA training and GPU inference benchmarks are therefore
planned on a free Colab or Kaggle GPU; measured hardware, runtime, precision,
sequence-length settings, and repetitions must be recorded when those experiments
are run.

\section{Results and Discussion}
\textbf{Pending execution.} Populate this section only with outputs produced by the
training, validation, final test, and benchmark scripts. Do not report the CPU
attention smoke test as a GPU result. Required final tables include baseline versus
fine-tuned classification, evidence retrieval, cache versus uncached latency and
memory, attention numerical error and runtime, and speculative acceptance and
throughput.

\section{Limitations}
The benchmark has 53 patient groups, and the selected 1.5B model is substantially
smaller than the models often used for complex clinical reasoning. Truncation can
remove relevant record sentences; the pipeline records this event but cannot recover
information not presented to the model. The output is a research prototype and is
not suitable for clinical decisions.

\section{Conclusion}
This draft specifies a reproducible criterion-level system and its evaluation
protocol. Conclusions about model quality, latency, and memory remain pending the
required CUDA experiments and untouched final test evaluation.

\section*{AI Use Statement}
AI assistance was used during implementation and drafting. The student must verify,
understand, and acknowledge all generated or adapted material and must independently
run the experiments and defend the final system.

\bibliographystyle{IEEEtran}
\bibliography{../references}
\end{document}
